% answers_short.tex —— 2002 年·春季京蒙皖卷·理科 · 简短答案速查
% 由 tools/make_short_answers.py 自 answers.tex 的 \gksolhead/\gkanswer 提取，
% 并与 2002/2002_数学真题_春季京蒙皖理_答案要点.md 逐题比对后生成，勿手改（改源 MD 后重跑）。

\begin{gknotice}
\textbf{使用说明}：本节只给最终答案便于快速核对，完整推导见解析版。答案为编写者独立推导，\textbf{非官方标准答案}。
\end{gknotice}

\gkpart{选择题}
\begin{center}\begin{tabular}{ccccccccccccc}
\toprule
题号 & 1 & 2 & 3 & 4 & 5 & 6 & 7 & 8 & 9 & 10 & 11 & 12 \\
\midrule
答案 & C & D & A & B & B & B & A & C & A & C & A & C \\
\bottomrule
\end{tabular}\end{center}

\gkpart{填空题}
\begin{itemize}[leftmargin=2.6em]
  \item[\textbf{13.}] $(\pm\sqrt7,0)$。
  \item[\textbf{14.}] $\dfrac{17\sqrt2}{26}$。
  \item[\textbf{15.}] $\dfrac{\sqrt2}{2}$。
  \item[\textbf{16.}] $90^\circ$。
\end{itemize}

\gkpart{解答题}
\begin{itemize}[leftmargin=2.6em]
  \item[\textbf{17.}] $\sqrt3$。
  \item[\textbf{18.}] 在 $(-\infty,0)$ 上为减函数。
  \item[\textbf{19.}] (1) 见解析；(2) $\dfrac\pi3$；(3) $\arccos\dfrac1{\sqrt{17}}$
  \item[\textbf{20.}] (1) 平均每年至少下降 $2$ 万元；(2) 一定够买
  \item[\textbf{21.}] (1) $x_n=\dfrac{x_{n-1}+x_{n-2}}2$；(2) $a_1=a$，$a_2=-\dfrac a2$，$a_3=\dfrac a4$，$a_n=a\left(-\dfrac12\right)^{n-1}$；(3) $\displaystyle\lim_{n\to\infty}x_n=\dfrac{2a}{3}$
  \item[\textbf{22.}] (1) $\dfrac{x^2}{25}+\dfrac{y^2}{9}=1$；(2) 弦 $AC$ 中点的横坐标为 $4$；(3) $-\dfrac{16}{5}<m<\dfrac{16}{5}$
\end{itemize}
